% Fragmento para incluir en la memoria del TFM.
% Requiere \usepackage{hyperref} en el preámbulo del documento principal.

\chapter{Repositorio de código y materiales reproducibles}
\label{anexo:repositorio-github}

El código, los notebooks, los resultados curados y la aplicación interactiva
desarrollados para este Trabajo Fin de Máster se encuentran disponibles en el
siguiente repositorio público de GitHub:

\begin{center}
    \href{https://github.com/xabibecerra/TFM---Xabier}%
    {\texttt{github.com/xabibecerra/TFM---Xabier}}
\end{center}

La aplicación desplegada puede consultarse en:

\begin{center}
    \href{https://tfm---xabier-cifshluvqdsrae5yvby4df.streamlit.app/}%
    {\texttt{tfm---xabier-cifshluvqdsrae5yvby4df.streamlit.app}}
\end{center}

El repositorio constituye el anexo técnico del trabajo e incluye la secuencia
de notebooks empleada en la preparación de datos, modelización, validación
temporal, interpretabilidad, análisis de errores y elaboración de candidatas
operativas. También contiene la aplicación desarrollada con Streamlit, los
scripts auxiliares, las pruebas automatizadas, los modelos finales congelados y
los archivos derivados necesarios para reproducir la demostración histórica.

La organización temporal publicada coincide con la utilizada en la memoria:
el entrenamiento emplea 2019, 2021 y enero--septiembre de 2022; octubre de 2022
se reserva para validación externa; y noviembre--diciembre de 2022 se utiliza
para una evaluación temporal retrospectiva con los modelos y preprocesadores
previamente congelados. La consulta exploratoria anterior de este último periodo
limita su independencia. Los viajes de enero--febrero de 2023 se presentan
exclusivamente como análisis descriptivo, ya que no se dispone de estados de
las estaciones para ese periodo.

Las copias públicas de los notebooks no contienen resultados de ejecución ni
rutas personales. Los resultados relevantes se facilitan en las carpetas
\texttt{results/} y \texttt{app\_data/}. Las importancias por ganancia y SHAP
deben interpretarse como señales utilizadas por el modelo y no como relaciones
causales. Del mismo modo, las propuestas operativas publicadas son candidatas
técnicas sujetas a restricciones logísticas y no órdenes de actuación reales.

Para favorecer la trazabilidad, el repositorio incorpora un archivo
\texttt{CITATION.cff}, huellas SHA-256 de los modelos y un flujo de integración
continua que ejecuta automáticamente las pruebas de la aplicación. La versión
de entrega queda identificada mediante la \emph{release} \texttt{v1.0.0-tfm},
publicada y consultada el 14 de septiembre de 2026. El código original se ofrece
bajo licencia MIT; las condiciones aplicables a los datos y materiales de
terceros se documentan por separado en \texttt{NOTICE.md}.
